% ============================================================================
% APÉNDICE A: CÓDIGO FUENTE CERTIFICADO EN SILICIO
% ============================================================================
\chapter{Código Fuente Certificado — POLYDIM V762/V764}


\section*{Nota sobre los Apéndices de Código}

El código fuente completo de la suite POLYDIM está disponible en:
\begin{itemize}
\item \textbf{Repositorio:} \url{https://github.com/AGT1973/POLYDIM_CLA_V7}
\item \textbf{Commit V762 certificado:} \texttt{f9fa786}
\item \textbf{Directorio local V762:} \texttt{E:\textbackslash POLYDIM\_EINSOF\textbackslash ENTREGA\_2026\_09\_19\_V762\textbackslash}
\item \textbf{Directorio local V764:} \texttt{E:\textbackslash POLYDIM\_EINSOF\textbackslash EPOLYDIM\_EINSOFENTREGA\_2026\_09\_20\_V764\textbackslash}
\end{itemize}

Los archivos de código fuente referenciados en esta tesis son:
\begin{tabular}{llll}
\hline
\textbf{Archivo} & \textbf{Versión} & \textbf{Descripción} \\
\hline
\texttt{kernel\_cpp\_v762.cpp} & V762 & Kernel C++20: Rodrigues, PMTP, Cayley-SMW \\
\texttt{kernel\_rust\_v762.rs} & V762 & Guardián Rust: Higham bound, Betti-1 DSU \\
\texttt{polydim\_v762\_monolito.py} & V762 & Orquestador Python con HardwareProbe \\
\texttt{polydim\_triton\_kernel\_v762.py} & V762 & Kernel Triton GPU FP64 \\
\texttt{test\_v762\_mpeleides.py} & V762 & Suite de 5 pruebas adversariales \\
\texttt{kernel\_cpp\_v764.cpp} & V764 & Kernel endurecido con 12 hallazgos A1--A12 \\
\texttt{kernel\_rust\_v764.rs} & V764 & Guardián Rust con tolerancias recalibradas \\
\texttt{polydim\_ffi\_v764.dart} & V764 & Dart FFI con medición real y liberación de memoria \\
\texttt{DOCUMENTO\_TESIS\_POLYDIM\_v764.md} & V764 & Documento fuente de esta tesis \\
\hline


\end{tabular}

\section{Fragmento: Neumaier Merge (V764)}


V764 introduce el método \texttt{merge} en el acumulador de Neumaier, que permite combinar
acumuladores de distintos hilos OpenMP de forma compensada:

\begin{verbatim}
struct Neumaier {
    double sum = 0.0;
    double c   = 0.0;
    inline void add(double v) {
        const double t = sum + v;
        if (std::abs(sum) >= std::abs(v)) c += (sum - t) + v;
        else                              c += (v   - t) + sum;
        sum = t;
    }
    // V764: merge dos acumuladores de forma compensada
    // Crítico: add(o.sum) + add(o.c) es más correcto que 
    // sum += o.sum; c += o.c (el último introduce un redondeo extra)
    inline void merge(const Neumaier& o) { add(o.sum); add(o.c); }
    inline double total() const { return sum + c; }
};
\end{verbatim}

\section{Fragmento: polydim\_orthonormalize\_pair\_f64 (V764)}


Función nueva en V764 que implementa Gram-Schmidt modificado de dos pasadas:

\begin{verbatim}
/* Dos pasadas de Gram-Schmidt modificado: una sola pasada deja un residuo
 * O(kappa*eps) que puede exceder la cota de 64*eps. */
for (int pass = 0; pass < 2; ++pass) {
    const double d = polydim_cdot(u, v, D, nt); // producto escalar compensado
    #pragma omp parallel for num_threads(nt) schedule(static)
    for (int64_t i = 0; i < static_cast<int64_t>(D); ++i)
        v[i] -= d * u[i];
}
\end{verbatim}

\section{Fragmento: polydim\_selftest\_compensation (V764)}


\begin{verbatim}
extern "C" POLYDIM_EXPORT int32_t POLYDIM_CALL polydim_selftest_compensation(void) {
    /* Si el compilador reasoció (a+b)+c == a+(b+c), la sustracción (s-t)
     * en Neumaier se vuelve algebraicamente cero y el compilador la elimina. */
    const double big = 1e15, small = 1.0;
    Neumaier acc;
    acc.add(big); acc.add(-big); acc.add(small);
    /* Con reasociación: acc.total() puede ser 0 en lugar de 1.0 */
    if (std::abs(acc.total() - small) > 0.5 * std::numeric_limits<double>::epsilon())
        return POLYDIM_ERR_COMPENSATION_BROKEN;
    return POLYDIM_SUCCESS;
}
\end{verbatim}
